\documentclass[11pt,a4paper]{article}

% ---------- Packages ----------
\usepackage[margin=1in]{geometry}
\usepackage{amsmath, amssymb, amsthm, mathtools}
\usepackage{microtype}
\usepackage{algorithm}
\usepackage{algpseudocode}
\usepackage[colorlinks=true,
            linkcolor=blue!50!black,
            urlcolor=blue!50!black,
            citecolor=blue!50!black]{hyperref}

% ---------- Theorem environments ----------
\theoremstyle{plain}
\newtheorem{theorem}{Theorem}[section]
\newtheorem{proposition}[theorem]{Proposition}
\newtheorem{lemma}[theorem]{Lemma}
\newtheorem{corollary}[theorem]{Corollary}

\theoremstyle{definition}
\newtheorem{definition}[theorem]{Definition}

\theoremstyle{remark}
\newtheorem{remark}[theorem]{Remark}
\newtheorem{example}[theorem]{Example}

% ---------- Math shortcuts ----------
\newcommand{\R}{\mathbb{R}}
\newcommand{\N}{\mathbb{N}}
\newcommand{\E}{\mathbb{E}}
\newcommand{\Var}{\operatorname{Var}}
\newcommand{\Cov}{\operatorname{Cov}}
\newcommand{\Corr}{\operatorname{Corr}}
\newcommand{\F}{\mathcal{F}}
\newcommand{\Prob}{\mathbb{P}}
\newcommand{\Q}{\mathbb{Q}}
\newcommand{\dd}{\mathop{}\!\mathrm{d}}
\newcommand{\eqdef}{\overset{\mathrm{def}}{=}}
\newcommand{\indic}{\mathbf{1}}

% ---------- Document metadata ----------
\title{Monte Carlo Pricing of a European Call:\\
       Antithetic Variates}
\author{Quant Finance Portfolio --- Phase 2, Block 2, Algorithm 1}
\date{\today}

\begin{document}
\maketitle

\begin{abstract}
We apply the antithetic variates technique of the Block~2.0 writeup
to the Monte Carlo pricer of the European call under the exact GBM
sampler of Block~1.1. The monotonicity condition holds trivially, so
the AV estimator is guaranteed to reduce variance over the IID
estimator at equal sample size, and to reduce VxT over the IID
estimator at equal computational budget. We predict the variance
reduction factor analytically as a function of moneyness ($S_0/K$)
and design the empirical validation that measures the factor and the
VxT ratio.
\end{abstract}

\tableofcontents

\section{Introduction}
\label{sec:intro}

Block~1.1 produced our baseline European-call pricer:
\begin{equation}
\hat C_M^{\mathrm{IID}} \;=\; \frac{1}{M} \sum_{i=1}^M
e^{-rT} (S_T^{(i)} - K)^+,
\quad
S_T^{(i)} = S_0 \exp\!\bigl((r - \tfrac{1}{2}\sigma^2)T + \sigma \sqrt T\, Z_i\bigr),
\end{equation}
with $Z_i \sim \mathcal{N}(0, 1)$ i.i.d. The empirical variance at
canonical parameters ($S_0 = K = 100$, $r = 0.05$, $\sigma = 0.20$,
$T = 1$) is $\sigma^2_Y \approx 217$, giving a half-width of about
$0.09$ at $M = 10^5$.

In this block we replace this estimator with its antithetic-variates
counterpart, derive the resulting variance under GBM, predict the
reduction factor as a function of moneyness, and validate empirically.

\section{The antithetic estimator for the European call}
\label{sec:av-call}

\subsection{Construction}

Let $f(z) \eqdef e^{-rT} (S(z) - K)^+$ where
\begin{equation}
S(z) \;\eqdef\; S_0 \exp\!\bigl((r - \tfrac{1}{2}\sigma^2)T + \sigma \sqrt T \, z\bigr),
\qquad z \in \R.
\end{equation}
The antithetic-pair payoff is
\begin{equation}
\label{eq:av-payoff}
Y_i^{\mathrm{AV}} \;\eqdef\; \tfrac{1}{2}\bigl(\, f(Z_i) + f(-Z_i) \,\bigr),
\end{equation}
with $Z_i \sim \mathcal{N}(0,1)$ i.i.d. The estimator is
\begin{equation}
\hat C_M^{\mathrm{AV}} \;\eqdef\; \frac{1}{M} \sum_{i=1}^M Y_i^{\mathrm{AV}}.
\end{equation}
Each $Y_i^{\mathrm{AV}}$ requires \emph{two} payoff evaluations (one
at $Z_i$, one at $-Z_i$), so the per-pair computational cost is
roughly twice that of an IID payoff. Per the VxT criterion of Block~2.0,
the relevant comparison is between $\hat C_M^{\mathrm{AV}}$ (using $2M$
payoff evaluations) and $\hat C_{2M}^{\mathrm{IID}}$ (also using $2M$
payoff evaluations).

\subsection{Verification of monotonicity}

The function $z \mapsto S(z)$ is strictly increasing in $z$ (it is
$S_0$ times a strictly increasing exponential). The function $s \mapsto
e^{-rT}(s - K)^+$ is monotone non-decreasing in $s$, strictly
increasing on $\{s > K\}$. The composition $z \mapsto f(z)$ is
therefore monotone non-decreasing, strictly increasing on $\{z : S(z)
> K\}$.

By Theorem~3.2 of the Block~2.0 writeup, $\Cov(f(Z), f(-Z)) \leq 0$
with strict inequality (since $f$ is not essentially constant: the
ITM and OTM regions both have positive probability under the standard
normal). The antithetic estimator is therefore strictly superior to
the IID estimator at equal sample size.

\subsection{Variance and reduction factor}

By Theorem~3.1 of the Block~2.0 writeup,
\begin{equation}
\label{eq:av-variance-call}
\Var(\hat C_M^{\mathrm{AV}}) \;=\; \frac{1}{2M}\bigl(\, \sigma^2_f + \Cov(f(Z), f(-Z)) \,\bigr),
\qquad \sigma^2_f \;\eqdef\; \Var(f(Z)).
\end{equation}
At equal computational budget,
\begin{equation}
\Var(\hat C_{2M}^{\mathrm{IID}}) \;=\; \frac{\sigma^2_f}{2M},
\end{equation}
so the variance reduction factor relative to IID-at-equal-cost is
\begin{equation}
\label{eq:vrf}
\mathrm{VRF} \;\eqdef\;
\frac{\Var(\hat C_{2M}^{\mathrm{IID}})}{\Var(\hat C_M^{\mathrm{AV}})}
\;=\; \frac{\sigma^2_f}{\sigma^2_f + \Cov(f(Z), f(-Z))}
\;=\; \frac{1}{1 + \Corr(f(Z), f(-Z))}.
\end{equation}
The $\mathrm{VRF}$ is exactly $1$ if $f(Z)$ and $f(-Z)$ are
uncorrelated (no benefit), goes to $\infty$ as the correlation
approaches $-1$ (perfect negative correlation, infinite reduction),
and equals $1/2$ if the correlation approaches $+1$ (AV is twice as
bad as IID --- but this cannot happen here, by the monotonicity
result).

In particular, $\mathrm{VRF} > 1$ if and only if
$\Corr(f(Z), f(-Z)) < 0$, which is exactly what monotonicity
guarantees.

\section{Predicting the VRF as a function of moneyness}
\label{sec:vrf-prediction}

How large is the correlation $\Corr(f(Z), f(-Z))$ for the European
call under GBM? The answer depends on moneyness, and the dependence
is informative.

\subsection{Two limiting regimes}

\paragraph{Deep ITM ($K \to 0$):} the call payoff approaches the
forward, $f(z) \to e^{-rT} S(z) - e^{-rT} K \cdot e^{rT} = e^{-rT}
S(z) - K e^{-rT}$. (For ATM and ITM cases below the strike $K$
appears as a constant after discounting.) The dominant variation is
in $S(z)$, which is monotone in $z$ in a strict, smooth way. The
function is essentially $z \mapsto a \cdot e^{bz}$ for constants
$a, b$. We expect $f(Z)$ and $f(-Z)$ to have strong negative
correlation (near $-1$), giving very high VRF.

\paragraph{Deep OTM ($K \to \infty$):} the call payoff is zero on
$\{S(z) < K\}$, which is the overwhelming majority of paths. Both
$f(Z)$ and $f(-Z)$ are zero with high probability, and one is
nonzero only when the other is zero (in opposite directions of $z$).
The correlation is small in absolute terms (close to zero) because
both samples are zero most of the time, giving VRF close to~$1$.

\paragraph{ATM ($K = S_0$):} the regime of practical interest. The
call activates roughly half the time, with the corresponding negative
correlation under reflection. We expect a moderate negative
correlation, around $-0.5$ to $-0.7$.

\subsection{The kink and the loss of correlation}

The reason the call payoff has \emph{less than perfect} antithetic
correlation, even though it is monotone, is the kink at $S = K$. On
the ITM side ($z$ large positive), $f(z) = e^{-rT}(S(z) - K)$
varies smoothly and approximately linearly in $S(z)$. On the OTM
side ($z$ large negative), $f(z) = 0$, identically. The reflection
$z \mapsto -z$ maps the ITM side to the OTM side and vice versa, but
the OTM side carries no information (zero), so the negative
correlation cannot be perfect.

In contrast, an option whose payoff is smooth and strictly monotone
everywhere (such as a plain forward, $f(z) = S(z) - F$) achieves
$\Corr(f(Z), f(-Z))$ very close to $-1$. The kink in the call payoff
caps the achievable correlation strictly above $-1$.

For our running ATM example we predict VRF in the range
$\mathrm{VRF} \approx 1.5$--$3.0$, which corresponds to negative
correlations in the range $-0.7$ to $-0.3$. The exact number will
emerge from the validation experiment of Section~\ref{sec:validation}.

\section{Algorithmic specification}
\label{sec:algorithm}

\begin{algorithm}[H]
\caption{Antithetic-variates Monte Carlo for the European call (exact GBM)}
\label{alg:mc-av}
\begin{algorithmic}[1]
\Require Spot $S_0$, strike $K$, maturity $T$, rate $r$, volatility
$\sigma$, sample size $M$, confidence level $1 - \beta$, random
generator \texttt{rng}.
\Ensure Estimate $\hat C^{\mathrm{AV}}$, half-width $h_M^{\mathrm{AV}}$,
sample variance $S_M^{2,\mathrm{AV}}$.
\State $D \gets e^{-rT}$,
$\quad \mu \gets (r - \tfrac{1}{2} \sigma^2) T$,
$\quad \nu \gets \sigma \sqrt T$
\For{$i = 1, \ldots, M$}
\State Draw $Z_i \sim \mathcal{N}(0, 1)$ using \texttt{rng}
\State $S_+^{(i)} \gets S_0 \exp(\mu + \nu Z_i)$
\State $S_-^{(i)} \gets S_0 \exp(\mu - \nu Z_i)$
\State $Y_i^{\mathrm{AV}} \gets \tfrac{1}{2}\,D\,
\bigl(\max(S_+^{(i)} - K, 0) + \max(S_-^{(i)} - K, 0)\bigr)$
\EndFor
\State $\hat C^{\mathrm{AV}} \gets M^{-1} \sum_i Y_i^{\mathrm{AV}}$
\State $S_M^{2,\mathrm{AV}} \gets (M-1)^{-1} \sum_i (Y_i^{\mathrm{AV}}
- \hat C^{\mathrm{AV}})^2$
\State $h_M^{\mathrm{AV}} \gets z_{\beta/2}\,
\sqrt{S_M^{2,\mathrm{AV}}/M}$
\State \Return $(\hat C^{\mathrm{AV}},\, h_M^{\mathrm{AV}},\, S_M^{2,\mathrm{AV}})$
\end{algorithmic}
\end{algorithm}

The algorithm uses the same standard normal $Z_i$ to drive both
elements of the antithetic pair, exploiting that $S(\nu Z) = S_0
e^{\mu + \nu Z}$ and $S(-\nu Z) = S_0 e^{\mu - \nu Z}$ require only
the sign flip in the exponent. No additional random draws are
needed beyond the $M$ standard normals already used by the IID
estimator.

The variance and half-width are computed from the \emph{paired
payoffs} $Y_i^{\mathrm{AV}}$, not from the individual $f(Z_i)$ and
$f(-Z_i)$. This is essential: the $Y_i^{\mathrm{AV}}$ are i.i.d.\ in
$i$ (each pair is independent of the rest), and $\mathtt{mc\_estimator}$
applied to them returns the correct variance and half-width
automatically. Treating the $2M$ individual payoff evaluations as
i.i.d.\ would be wrong: $f(Z_i)$ and $f(-Z_i)$ are highly correlated,
and the independence assumption would inflate the apparent half-width.

This is structurally important: \emph{the variance-reducing
correlation lives inside each pair, not across the dataset}, which
is what makes the AV estimator both correct and efficient.

\section{Empirical validation strategy}
\label{sec:validation}

The validation script for the AV pricer comprises four tests.

\paragraph{Test 1: variance reduction factor.} Estimate the empirical
correlation $\hat\rho \eqdef \widehat{\Corr}(f(Z), f(-Z))$ from a
large sample of $M$ standard normals, and compute the VRF
\begin{equation}
\widehat{\mathrm{VRF}} \;=\; \frac{1}{1 + \hat\rho}.
\end{equation}
For our canonical ATM call, the prediction is $\hat\rho \in [-0.7,
-0.3]$, $\widehat{\mathrm{VRF}} \in [1.4, 3.3]$. The test passes if
the empirical correlation is negative and the VRF is meaningfully
above $1$.

\paragraph{Test 2: consistency with the IID estimator.} The AV
estimator and the IID estimator should agree on the same path data
(when computed from the same RNG state, which they will in our
design). Concretely: at $M = 10^4$ pairs, run AV pricer and IID
pricer with $M_{\text{IID}} = 2M$ paths from the same seed, and
verify that both estimates fall within a few half-widths of the
closed-form Black-Scholes price. (We do not require pathwise
equality of the two estimates because the path constructions
differ: the AV pricer uses each $Z_i$ for both elements of the
pair, while the IID pricer would treat them as $2M$ independent
draws.)

\paragraph{Test 3: VxT comparison at fixed budget.} Run both
estimators with the same number of payoff evaluations ($2M$ each)
and report:
\begin{itemize}
\item Empirical sample variance of $Y^{\mathrm{AV}}$ vs sample
variance of $Y^{\mathrm{IID}}$ at the same accumulated number of
payoffs.
\item Half-widths of the two estimators.
\item Ratio (which should equal $\widehat{\mathrm{VRF}}$ from Test~1
up to noise).
\end{itemize}
Wall-clock measurements are not necessary for the test to pass; the
per-pair compute is dominated by two payoff evaluations and one
shared random draw, with negligible overhead, so the VxT ratio
equals the variance ratio to within rounding.

\paragraph{Test 4: input validation.} Mechanical: the AV pricer
rejects bad inputs and accepts negative interest rates.

\section{Bibliographic notes}
\label{sec:biblio}

The textbook reference is Glasserman~\cite[Section~4.2]{Glasserman}.
The general theorem on antithetic variates and monotone functions
goes back to Hammersley \& Morton; the application to financial Monte
Carlo is in Boyle, Broadie \& Glasserman~\cite{BBG1997}.

\begin{thebibliography}{99}

\bibitem{BBG1997}
P.\ Boyle, M.\ Broadie, and P.\ Glasserman,
\emph{Monte Carlo methods for security pricing},
Journal of Economic Dynamics and Control, 21(8--9):1267--1321, 1997.

\bibitem{Glasserman}
P.\ Glasserman,
\emph{Monte Carlo Methods in Financial Engineering},
Springer, 2004.

\end{thebibliography}

\end{document}
